\section{Results}
\label{sec:results}

\subsection{A 2.7M-parameter model against foundation models}
\label{sec:main}

Table~\ref{tab:main} compares the axis-free CroFreMo against the strongest
reproduced baseline on each corpus and against our own tri-axial encoder;
the full matrix of fifteen baselines is in the appendix. The axis-free
model wins on the two seizure-detection corpora with adequate positives
(TUSZ, CHB-MIT) by margins far outside the baselines' seed variance, wins
narrowly on IIIC and TUEP, ties the strongest foundation model on TUEV and
the strongest supervised model on TUAB, and loses on the two sleep-staging
corpora, on the two cohort-diagnosis corpora, on TUAR and on Siena. On
CHB-MIT the strongest baseline is not the largest model: CBraMod trained
from scratch reaches $0.317 \pm 0.167$ and its released pretrained weights
$0.233 \pm 0.080$, against $0.716$ for a model thirteen times smaller.
Against the tri-axial encoder it was derived from, the axis-free model
gains $0.06$ on TUSZ and $0.02$ on CHB-MIT, matches it within a point on
six corpora, and is $0.02$--$0.04$ lower on IIIC, TUAR and ADFD.

\begin{table}[t]
\centering
\small
\caption{CroFreMo against the strongest reproduced baseline per corpus.
Baselines and the tri-axial CroFreMo are three-seed means $\pm$ std unless
marked $(1)$; the axis-free CroFreMo (2.7M, the final model) is a single
seed throughout. ``pre'' marks a baseline evaluated with its released
pretrained weights. TUAB is scored by balanced accuracy for every model.}
\label{tab:main}
\resizebox{\textwidth}{!}{%
\begin{tabular}{lllllc}
\toprule
Corpus & Metric & Strongest baseline & & CroFreMo tri-axial (1.6M) & CroFreMo axis-free (2.7M) \\
\midrule
TUSZ & AUC-PR & FFCL & $0.545 \pm 0.024$ & $0.639 \pm 0.029$ & $\mathbf{0.698}\,(1)$ \\
CHB-MIT & AUC-PR & TFM-Tokenizer pre & $0.627 \pm 0.021$ & $0.698 \pm 0.045$ & $\mathbf{0.716}\,(1)$ \\
IIIC & $\kappa$ & REVE pre & $0.436 \pm 0.002$ & $\mathbf{0.479 \pm 0.008}$ & $0.447\,(1)$ \\
TUEP & AUROC & EEGPT scratch & $0.786 \pm 0.018$ & $\mathbf{0.810 \pm 0.003}$ & $0.799\,(1)$ \\
TUEV & $\kappa$ & REVE pre & $0.685 \pm 0.032$ & $0.690 \pm 0.029$ & $0.680\,(1)$ \\
TUAB & bal.\ acc. & ST-Transformer & $0.820 \pm 0.004$ & $0.816\,(1)$ & $0.820\,(1)$ \\
Sleep-EDF & $\kappa$ & ContraWR & $0.692 \pm 0.012$ & $0.642 \pm 0.011$ & $0.639\,(1)$ \\
CAUEEG & bal.\ acc. & BIOT scratch & $0.561 \pm 0.009$ & $0.525 \pm 0.012$ & $0.511\,(1)$ \\
ISRUC & $\kappa$ & CBraMod pre & $0.754 \pm 0.006$ & $0.702 \pm 0.007$ & $0.698\,(1)$ \\
ADFD & bal.\ acc. & BIOT scratch & $0.525 \pm 0.017$ & $0.505 \pm 0.050$ & $0.461\,(1)$ \\
TUAR & $\kappa$ & CBraMod pre & $0.715\,(1)$ & $0.620 \pm 0.030$ & $0.600\,(1)$ \\
Siena & AUC-PR & REVE pre & $0.518 \pm 0.096$ & $0.170 \pm 0.070$ & $0.201\,(1)$ \\
\bottomrule
\end{tabular}}
\end{table}

\subsection{What the seizure margin is made of}
\label{sec:attribution}

Table~\ref{tab:attribution} is the experiment the encoder was built for.
The axis-free model is trained with its tokenizer producing either the
duplex grid of Section~\ref{sec:tokenizer} or waveform tokens only, with
every other setting identical, at the adopted width and at one width
larger. Removing the coupling content costs $0.136$ AUC-PR on CHB-MIT,
$0.107$ $\kappa$ on TUEV, $0.063$ AUC-PR on TUSZ and $0.027$ $\kappa$ on
IIIC at $d{=}192$. At $d{=}256$ the four effects keep their sign; three
are of similar size and the CHB-MIT effect drops to $0.025$, so the size
of the effect is sensitive to capacity and optimization state, not only
to the corpus. For scale, the seed standard deviations of the tri-axial
model on the same corpora are $0.029$, $0.045$, $0.029$ and $0.008$; they
are a reference for the recipe, not a substitute for the axis-free
model's own variance. The two margins by which the model leads the
foundation-model field are therefore not properties of the band
decomposition or of the folded attention alone: without the interaction
token the same encoder falls to $0.635$ on TUSZ and $0.580$ on CHB-MIT,
below the tri-axial model and, on CHB-MIT, below TFM-Tokenizer.

This comparison changes more than one thing at once --- the waveform-only
grid has half the rows, no explicit envelope feature and no phase
representation --- so it attributes the gain to the interaction rows as a
whole, not yet to the cross-frequency alignment inside them. Controls
that hold the row count fixed and vary what the extra rows carry are the
subject of Section~\ref{sec:design}.

\begin{table}[t]
\centering
\small
\caption{Coupling attribution in the axis-free encoder: duplex grid vs.\
waveform tokens only, everything else identical. Single seeds.}
\label{tab:attribution}
\begin{tabular}{llcccccc}
\toprule
 & & \multicolumn{3}{c}{$d{=}192$ (final)} & \multicolumn{3}{c}{$d{=}256$} \\
\cmidrule(lr){3-5}\cmidrule(lr){6-8}
Corpus & Metric & coupling on & off & $\Delta$ & on & off & $\Delta$ \\
\midrule
CHB-MIT & AUC-PR & $\mathbf{0.716}$ & $0.580$ & $\mathbf{+0.136}$ & $0.608$ & $0.583$ & $+0.025$ \\
TUEV & $\kappa$ & $\mathbf{0.680}$ & $0.573$ & $\mathbf{+0.107}$ & $0.679$ & $0.601$ & $+0.078$ \\
TUSZ & AUC-PR & $\mathbf{0.698}$ & $0.635$ & $\mathbf{+0.063}$ & $0.631$ & $0.541$ & $+0.090$ \\
IIIC & $\kappa$ & $\mathbf{0.447}$ & $0.420$ & $+0.027$ & $0.465$ & $0.418$ & $+0.047$ \\
\bottomrule
\end{tabular}
\end{table}

\subsection{The same token in four encoders}
\label{sec:content-structure}

Table~\ref{tab:content-structure} places the same coupling content inside
four encoders. The tri-axial encoder, which attends along a dedicated
frequency axis, gains $0.15$ from it on TUEV and nothing outside seed
variance on the other corpora (the full three-seed table is in the
appendix); its seizure-detection margin over the field is carried by the
band decomposition and the factorized attention (Table~\ref{tab:arch}).
CBraMod's encoder, which keeps its own rFFT magnitude branch, gains from
the interaction rows relative to matched waveform rows on TUEV ($0.045$),
TUEP ($0.057$), IIIC ($0.015$) and CHB-MIT ($0.124$), and loses $0.092$ on
TUSZ. LaBraM and REVE, trained from scratch with no spectral pathway,
gain $0.150$ and $0.241$ on TUEV when the tokenizer replaces their patch
embedding. The axis-free encoder gains on all four corpora measured.

\begin{table}[t]
\centering
\small
\caption{The same coupling content in four encoders: change in the primary
metric when it is switched on (tri-axial, axis-free: duplex vs.\ waveform
tokens; CBraMod: interaction rows vs.\ matched waveform rows added to the
native tokenizer; LaBraM, REVE: CroFreMo tokenizer vs.\ native, hosts from
scratch). Tri-axial: three seeds; the rest single seeds.}
\label{tab:content-structure}
\resizebox{\textwidth}{!}{%
\begin{tabular}{llccccc}
\toprule
Encoder & Own spectral pathway & TUEV & TUSZ & CHB-MIT & IIIC & TUEP \\
\midrule
Tri-axial CroFreMo & frequency attention axis & $+0.15$ & $-0.03$ & $+0.03$ & $+0.01$ & $0.00$ \\
CBraMod + rows & rFFT magnitude branch & $+0.045$ & $-0.092$ & $+0.124$ & $+0.015$ & $+0.057$ \\
LaBraM, tokenizer replaced & none & $+0.150$ & --- & $+0.017$ & $+0.014$ & --- \\
REVE, tokenizer replaced & none & $+0.241$ & --- & --- & $+0.078$ & --- \\
Axis-free CroFreMo & none & $+0.107$ & $+0.063$ & $+0.136$ & $+0.027$ & --- \\
\bottomrule
\end{tabular}}
\end{table}

The pattern is consistent with a simple reading, which we state as a
reading and not as a mechanism we have isolated. An encoder with a
dedicated frequency axis or a spectral branch has other routes to
band-power and, through attention over band tokens, to cross-band
relations; for it the explicit content is redundant except where the
relation is specific enough to be expensive to re-derive, as on
periodic epileptiform discharges, where the tri-axial gain concentrates
(GPED $+0.27$, PLED $+0.09$ in per-class F1,
Table~\ref{tab:perclass} in the appendix). An encoder without such a
route --- the folded-attention encoder, LaBraM, REVE --- can in
principle still learn cross-band relations from waveform rows through
its attention, but in the label regime of these corpora it benefits
substantially from being handed the relation. The negative TUSZ cell and
the positive CHB-MIT cell in CBraMod, on two seizure-detection corpora,
show that the corpus and the host's own training state matter as much
as the presence of a spectral pathway: CBraMod trains poorly on CHB-MIT
from scratch and the interaction rows give it a foothold there. Whether
an encoder with a spectral pathway actually computes the coupling that
the token supplies is a question for probing or intervention
experiments, not one this table answers.

\subsection{Design controls}
\label{sec:design}

Table~\ref{tab:design} reports the axis-free variants that were tried on
the way to the final model, all single seeds, alongside the tri-axial
architectural controls of Table~\ref{tab:arch}. Folding the band rows
into spatial attention rather than leaving them as unrelated channels is
worth $0.09$ on both seizure corpora and costs $0.02$--$0.04$ on the
smaller corpora, where the longer attention span overfits. Raising the
width from 192 to 256 recovers those small-corpus cells and loses
$0.07$ and $0.11$ on TUSZ and CHB-MIT. An explicit coupling-strength
feature --- the magnitudes $|Z_{ij}|$ projected onto the fused rows ---
raises TUEV to $0.721$ and Siena to $0.398$ and collapses CHB-MIT to
$0.513$. No variant dominates; the adopted one is the one that keeps the
seizure margin, and every alternative is reported. Within the tri-axial
family, collapsing the learned filterbank to a single broadband band
costs $0.11$ on TUSZ and $0.42$ on CHB-MIT, and replacing factorized
attention by a flat encoder that mixes electrodes before attention costs
$0.22$ and $0.53$: the band-decomposed grid and attention factorized over
it are what let a small model learn rare-positive seizure windows at all.

\begin{table}[t]
\centering
\small
\caption{Axis-free design variants, single seeds. ``fold'' = band rows
folded into spatial attention; ``strength'' = explicit $|Z|$ feature.}
\label{tab:design}
\resizebox{\textwidth}{!}{%
\begin{tabular}{lcccccccc}
\toprule
Variant & TUSZ & CHB-MIT & TUEV & IIIC & TUAR & ADFD & Siena & Sleep-EDF \\
\midrule
no fold, $d{=}192$ & 0.612 & 0.630 & 0.690 & 0.458 & 0.624 & 0.498 & 0.239 & 0.610 \\
no fold, $d{=}192$, strength & 0.643 & 0.513 & $\mathbf{0.721}$ & 0.448 & 0.626 & 0.467 & $\mathbf{0.398}$ & 0.614 \\
fold, $d{=}192$ (final) & $\mathbf{0.698}$ & $\mathbf{0.716}$ & 0.680 & 0.447 & 0.600 & 0.461 & 0.201 & 0.639 \\
fold, $d{=}256$ & 0.631 & 0.608 & 0.679 & $\mathbf{0.465}$ & $\mathbf{0.657}$ & 0.454 & 0.269 & $\mathbf{0.653}$ \\
\bottomrule
\end{tabular}}
\end{table}

\begin{table}[t]
\centering
\small
\caption{Tri-axial architectural controls, waveform-token frontend, three
seeds. Removing either the band decomposition or the factorization
collapses seizure detection.}
\label{tab:arch}
\begin{tabular}{lccc}
\toprule
 & 8 bands, tri-axial & 1 band, tri-axial & 8 bands, flat (0.9M) \\
\midrule
TUSZ (AUC-PR) & $0.671 \pm 0.026$ & $0.558 \pm 0.041$ & $0.448 \pm 0.048$ \\
CHB-MIT (AUC-PR) & $0.667 \pm 0.047$ & $0.245 \pm 0.068$ & $0.136 \pm 0.016$ \\
\bottomrule
\end{tabular}
\end{table}

\subsection{The token inside published encoders}
\label{sec:transplant}

Table~\ref{tab:transplant} adds the tokenizer's rows to CBraMod as
described in Section~\ref{sec:encoder}. Against the matched control ---
the same eight extra rows carrying waveform tokens --- the interaction
rows are better on TUEV, TUEP, IIIC and CHB-MIT and worse on TUSZ, so
the effect is attributable to what the rows carry rather than to the
extra channels. Against CBraMod's own tokenizer the transplant is a clear
gain on TUEV ($+0.077$, four standard deviations of the native model) and
on CHB-MIT ($+0.148$, where the native model is unstable across seeds), a
tie on TUEP and IIIC, and a loss on TUSZ.
Table~\ref{tab:hosts} widens the test to two further encoders, with
the tokenizer now \emph{replacing} the host's patch embedding (every
electrode--row pair becomes one host channel, rows are pooled after the
encoder, the head is unchanged). Trained from scratch, LaBraM gains
$0.150$ $\kappa$ on TUEV and REVE $0.241$; on IIIC the gains are
$0.014$ and $0.078$; on CHB-MIT LaBraM gains $0.017$, within its own
seed variance. The three hosts span 4M to 69M parameters and three
encoder designs, and the ordering of the TUEV gains follows how far each
host is from what it can learn from $10^5$ labelled windows on its own.
Adding the same rows to CBraMod's \emph{released} checkpoint is worse
than the checkpoint alone on both corpora tried (appendix): a pretrained
encoder is tuned to its own token distribution, and inserting eight new
channels per electrode disturbs it more than the coupling content
repays. We draw the conclusion the data support --- the token transfers
into encoders trained from scratch, its effect is attributable to the
content of the added rows, and it pays off on epileptiform event
classification --- and not the broad one about released checkpoints.

\begin{table}[t]
\centering
\small
\caption{CroFreMo's rows added to CBraMod's encoder, trained from scratch.
CBraMod's own tokenizer is a three-seed mean; the two transplant arms are
single seeds with identical parameter and token counts.}
\label{tab:transplant}
\begin{tabular}{llccc}
\toprule
Corpus & Metric & CBraMod tokenizer & + waveform rows & + interaction rows \\
\midrule
TUEV & $\kappa$ & $0.564 \pm 0.019$ & $0.596$ & $\mathbf{0.641}$ \\
CHB-MIT & AUC-PR & $0.317 \pm 0.167$ & $0.341$ & $\mathbf{0.465}$ \\
TUEP & AUROC & $0.642 \pm 0.028$ & $0.593$ & $\mathbf{0.650}$ \\
IIIC & $\kappa$ & $0.393 \pm 0.005$ & $0.381$ & $0.396$ \\
TUSZ & AUC-PR & $\mathbf{0.482 \pm 0.043}$ & $0.476$ & $0.384$ \\
\bottomrule
\end{tabular}
\end{table}

\begin{table}[t]
\centering
\small
\caption{CroFreMo's tokenizer in place of the host's own, hosts trained
from scratch. Native rows are three-seed means except REVE from scratch
(single seed); transplants are single seeds. Cohen's $\kappa$; CHB-MIT
AUC-PR.}
\label{tab:hosts}
\resizebox{\textwidth}{!}{%
\begin{tabular}{lcccccc}
\toprule
Host (from scratch) & TUEV native & TUEV + CroFreMo & IIIC native & IIIC + CroFreMo & CHB-MIT native & CHB-MIT + CroFreMo \\
\midrule
CBraMod (4M) & $0.564 \pm 0.019$ & $\mathbf{0.632}$ & $0.393 \pm 0.005$ & $0.396$ & --- & --- \\
LaBraM (5.9M) & $0.372 \pm 0.009$ & $\mathbf{0.522}$ & $0.406 \pm 0.007$ & $0.420$ & $0.360 \pm 0.050$ & $0.377$ \\
REVE (69M) & $0.303$ & $\mathbf{0.544}$ & $0.299$ & $\mathbf{0.377}$ & --- & --- \\
\bottomrule
\end{tabular}}
\end{table}

\subsection{Pretraining}
\label{sec:pretrain-results}

The axis-free model is trained from scratch throughout; it has no
pretrained row. The tri-axial model was pretrained with the objective of
Section~\ref{sec:objective} and finetuned under one fixed recipe, and the
outcome --- gains where labels are scarce or noisy, losses where they are
plentiful --- is reported as an exploratory result in the appendix
(Table~\ref{tab:pretrain-triaxial}), together with the comparison of
objectives. We do not build on it in the main text.
